\chapter{Robust Sequential Tests at Unbounded Endpoints}
\label{sec:p15-robust-sequential-endpoints}

An infinite relative entropy does not make an expected stopping time
asymptotic to zero. It changes the scale of the answer. This chapter
completes the expected-sample-size argument in Huber's Section 5,
including zero clipping endpoints and singular observations
\cite[p.~1757]{Huber1965RobustProbabilityRatio}.
The key comparison is with a one-sided hitting time; the proof does
not require a second moment or a bounded logarithmic likelihood ratio.

\begin{hbdefinition}{def:p15-sequential-endpoint-setting}{The unchanged neighborhoods and oriented scores}
Use a separated pair from \cref{sec:p15-robust-hypothesis-testing}.
On the same arbitrary measurable observation space, let
\(\mathcal H_i\) be either the original contamination classes
\[
 \mathcal H_i=\{(1-\varepsilon_i)P_i+\varepsilon_iH:
                    H\text{ any probability}\},\qquad 0\leq\varepsilon_i<1,
\]
or the original equal-radius variation classes
\[
 \mathcal H_i=\{R:2\sup_E|R(E)-P_i(E)|\leq\varepsilon\}.
\]
The classes are assumed disjoint, with their normalized least-favorable
members \(Q_i\) and their everywhere-defined clipped score \(\ell\)
from \cref{thm:p15-huber-contamination-pair,thm:p15-huber-tv-pair}.
The variation radius is nonnegative and uses exactly the convention
in \cref{def:p15-huber-neighborhoods}.
Zero radii and clipping endpoints zero or infinity are retained.
The ordered composite classes of \cref{thm:p15-robust-mlr-extension}
are also permitted, with its proved score order.

For each fixed \(R\in\mathcal H_i\), observations are independent
with common law \(R\). Write \(\gamma=\log\ell\), with extended values,
and let \(N_{a,b}\) be the first exit of \(\sum_{j=1}^n\gamma(X_j)\)
from \((-a,b)\), where \(a,b>0\). An infinite score causes immediate
exit; no subsequent opposite infinity is added.
The correct and wrong boundaries, in coordinates directed toward the
true hypothesis, are
\[
 Y_0=-\gamma,\quad Y_1=\gamma,\qquad
 (t_0,u_0)=(a,b),\quad(t_1,u_1)=(b,a).
\]
Let \(Z_i\) have the law of \(Y_i\) under \(Q_i\).
For a real or extended \(z\), set \(z^+=\max(z,0)\),
\(z^-=\max(-z,0)\). Define
\[
 D_i=E_{Q_i}Y_i,\qquad s_i=Q_i(Y_i=+\infty).
\]
The next result proves that these expectations are well defined and
that \(D_i\in(0,\infty]\).
\end{hbdefinition}

\begin{hblemma}{thm:p15-directional-log-integrability}{proved=annex:p15.directional.log.integrability}{Negative parts, entropy and singular mass}
For probabilities \(Q\ne P\), choose finite densities \(q,p\) with
respect to \(\nu=Q+P\). Define \(Z=\log(q/p)\) under \(Q\),
with value \(+\infty\) where \(q>0=p\). Then
assign any finite value on the \(Q\)-null set \(\{q=0\}\). We have
\[
 Q(Z=-\infty)=0,\qquad EZ^-\leq 1/e,\qquad
 E e^{-Z}=P(q>0)\leq1.
\]
Its extended expectation is the relative entropy
\(D(Q\Vert P)\in(0,\infty]\).
It is finite only if \(Q\ll P\); positive singular mass forces
infinite entropy, but infinite entropy can also occur with equivalent
laws and finite scores everywhere.

In \cref{def:p15-sequential-endpoint-setting}, every \(Y_i\) under
\(R\in\mathcal H_i\) is stochastically no smaller than \(Z_i\).
Consequently
\[
 E_RY_i^-\leq EZ_i^-\leq1/e,\qquad
 E_RY_i\geq D_i>0,\qquad E_Re^{-Y_i}\leq1.
\]
These assertions allow the positive expectation to be infinite.
\end{hblemma}

\begin{hbdefinition}{def:p15-one-sided-hitting}{A one-sided passage and its mean}
For independent copies of \(Z\in\mathbb R\cup\{+\infty\}\), put
\(S_0=0,S_n=\sum_{j=1}^nZ_j\). For \(t>0\), let
\[
 \tau_t=\inf\{n\geq1:S_n\geq t\},\qquad m(t)=E\tau_t,
 \quad \tau_0=0,\quad m(0)=0.
\]
The infimum of an empty set is infinity. A positive infinite increment
makes every later partial sum infinite. Define
\(d_M=E\min(Z,M)\) for a finite \(M>0\).
\end{hbdefinition}

\begin{hbtheorem}{thm:p15-one-sided-truncation}{proved=annex:p15.one.sided.truncation}{Finite passage means without an overshoot moment assumption}
Suppose \(EZ^-<\infty\) and \(EZ=D\in(0,\infty]\).
There is a finite \(M>0\) with \(d_M>0\). For every such \(M\)
and \(t>0\),
\[
 1\leq m(t)\leq\frac{t+M}{d_M}<\infty.
\]
If \(D<\infty\), the untruncated stopped sum is integrable and
\[
 ES_{\tau_t}=D\,m(t),\qquad m(t)\geq t/D,\qquad
 \frac{m(t)}t\longrightarrow\frac1D\quad(t\to\infty).
\]
If \(D=\infty\), the correct normalized statement is
\(m(t)/t\to0\). No expression \(m(t)\sim t/\infty\) is used.

If \(s=\Pr(Z=\infty)>0\), then \(m(t)\to1/s\).
If \(s=0\), then \(m(t)\to\infty\), including the case \(D=\infty\).
\end{hbtheorem}

\begin{hblemma}{thm:p15-hitting-restart}{proved=annex:p15.hitting.restart}{Measurable restarting and subadditivity}
Under the preceding passage assumptions, \(m\) is finite,
nondecreasing and Borel measurable on \([0,\infty)\), and
\[
 m(t+v)\leq m(t)+m(v)\qquad(t,v\geq0).
\]
Every finite future block after an almost surely finite stopping time
has the original product law and is independent of the stopped past.
In particular, if
\[
 N=\inf\{n\geq1:S_n\notin(-u,t)\},\qquad u,t>0,
\]
and \(W=\{S_N\leq-u\}\) is its wrong-exit event, then
\[
 m(t)=EN+E\bigl[{\bf1}_W m(t-S_N)\bigr].
\]
The right-hand threshold on \(W\) is finite and positive.
The annex proves the restarting identity on the real score space,
including its measurable random-threshold argument.
\end{hblemma}

\begin{hbtheorem}{thm:p15-two-boundary-passage-comparison}{proved=annex:p15.two.boundary.passage.comparison}{A quantitative comparison uniform in the correct boundary}
Under the assumptions of \cref{thm:p15-one-sided-truncation}, assume
also \(Ee^{-Z}\leq1\), and fix \(M>0\) with \(d_M>0\).
For the two-boundary stop above and \(u\geq1,t>0\), put
\[
 \epsilon_M(u)=e^{-u}\left(1+\frac{u+M}{d_M}\right).
\]
Then
\[
 E e^{-S_N}\leq1,\qquad \Pr(W)\leq e^{-u},\qquad
 1-\epsilon_M(u)\leq\frac{EN}{m(t)}\leq1.
\]
In particular \(EN/m(t)\to1\) as \(u\to\infty\), uniformly over
all \(t>0\). This conclusion does not require a bounded increment,
a finite variance, a finite positive log moment, or a relative
growth condition between \(t\) and \(u\).

There is also a bound using the actual wrong-exit probability \(p=\Pr(W)\),
valid for any \(t,u>0\). For every \(K\geq1\),
\[
 0\leq1-\frac{EN}{m(t)}
 \leq p+\frac{(K+M)(p+e^{-K})}{d_M}.
\]
If \(0<p\leq e^{-1}\), choose \(K=\log(1/p)\); the upper bound becomes
\[
 \delta_M(p)=p\left\{1+\frac{2(\log(1/p)+M)}{d_M}\right\}\longrightarrow0.
\]
At \(p=0\) the ratio is exactly one and set \(\delta_M(0)=0\).
Thus a vanishing wrong-exit probability suffices even if its boundary
does not diverge.
\end{hbtheorem}

\begin{hbtheorem}{thm:p15-robust-endpoint-closure}{proved=annex:p15.robust.endpoint.closure}{Worst-case expected size at every entropy endpoint}
Retain exactly the neighborhoods in
\cref{def:p15-sequential-endpoint-setting}. Let \(m_i(t)\) be the
one-sided passage mean for the least-favorable score \(Z_i\).
For each \(i\), choose finite \(M_i>0\) with
\(d_i=E\min(Z_i,M_i)>0\), and put
\(\epsilon_i(u)=e^{-u}\{1+(u+M_i)/d_i\}\).
For \(u_i\geq1\) and \(t_i>0\),
\[
 (1-\epsilon_i(u_i))m_i(t_i)
 \leq E_{Q_i}N_{a,b}
 \leq \sup_{R\in\mathcal H_i}E_RN_{a,b}
 \leq m_i(t_i).
\]
Every expectation in this display is finite.
Thus, as \(a,b\to\infty\) in any manner,
\[
 \frac{\sup_{R\in\mathcal H_i}E_RN_{a,b}}{E_{Q_i}N_{a,b}}
 \longrightarrow1 .
\]
The supremum is over the entire original neighborhood, with a fixed
common one-observation law in each expectation.
The ratio conclusion holds for finite or infinite \(D_i\) and
for \(s_i>0\) as well as \(s_i=0\).
More generally, along any sequence of positive boundary pairs, let
\(p_i=\Pr_{Q_i}(\text{wrong exit})\). If \(p_i\to0\), then
\[
 \frac{E_{Q_i}N_{a,b}}{m_i(t_i)}\longrightarrow1,\qquad
 \frac{\sup_{R\in\mathcal H_i}E_RN_{a,b}}{E_{Q_i}N_{a,b}}
                    \longrightarrow1,
\]
using the actual-error bound above. In particular both source
expected-size ratios hold whenever both least-favorable error
probabilities tend to zero; boundary divergence is sufficient,
not necessary.

For the following size classifications assume \(a,b\to\infty\).
Then both \(E_{Q_i}N_{a,b}/t_i\) and
\(\sup_R E_RN_{a,b}/t_i\) tend to \(1/D_i\) when \(D_i<\infty\),
and to zero when \(D_i=\infty\).
When \(s_i>0\), the unnormalized two expectations tend to \(1/s_i\).
When \(s_i=0\), they both tend to infinity, even if \(D_i=\infty\).
These results close the unbounded-score endpoint of the source's
asymptotic least-favorability claim.
For each fixed \(R\in\mathcal H_i\), the same directional formula holds
with \(D_i\) replaced by \(E_RY_i\), using a zero normalized limit
when that mean is infinite. If \(E_RY_i=D_i<\infty\), the two
directional score laws coincide and their stopping-time laws agree
at every pair of boundaries. Equality of two infinite means does
not imply that conclusion.
\end{hbtheorem}

\section{What optimality over sequential procedures means}
\label{sec:p15-sequential-cost-minimax}

\begin{hbdefinition}{def:p15-sequential-procedure-cost}{Admissible sequential decisions and minimax expected cost}
A competing sequential procedure uses measurable stopping and decision
rules on finite observation histories, optionally with an independent
random seed having the same specified law under every hypothesis.
Its stopping time \(T\geq1\) is almost surely finite under every allowed
law, and its terminal decision \(\Delta\in\{0,1\}\) is measurable at
stopping; decision one rejects the null.
For \(0<\alpha,\beta<1\), \(\alpha+\beta<1\), require
\[
 \sup_{R\in\mathcal H_0}\Pr_R(\Delta=1)\leq\alpha,\qquad
 \sup_{R\in\mathcal H_1}\Pr_R(\Delta=0)\leq\beta.
\]
For each \(i\), let
\[
 \mathcal M_i(\alpha,\beta)=
 \inf_{(T,\Delta)\ {\rm satisfying\ these\ constraints}}
                 \sup_{R\in\mathcal H_i}E_RT .
\]
Expectations may be infinite. The two infima are coordinatewise
cost criteria; a single procedure need not attain either infimum.
\end{hbdefinition}

\begin{hbtheorem}{thm:p15-stopped-binary-information}{proved=annex:p15.stopped.binary.information}{A complete information lower bound at stopping}
For two simple laws \(Q_0,Q_1\), let a procedure satisfy the displayed
error bounds at those laws and terminate almost surely at both.
If \(0<D_0=D(Q_0\Vert Q_1)<\infty\), then
\[
 E_{Q_0}T\geq
 \frac{(1-\alpha)\log((1-\alpha)/\beta)
                +\alpha\log(\alpha/(1-\beta))}{D_0}.
\]
If \(0<D_1<\infty\), then
\[
 E_{Q_1}T\geq
 \frac{(1-\beta)\log((1-\beta)/\alpha)
                +\beta\log(\beta/(1-\alpha))}{D_1}.
\]
The inequalities include infinite expected stopping times.
The proof constructs the stopped-history densities, proves the
stopped entropy identity, and derives the binary decision comparison.
\end{hbtheorem}

\begin{hbtheorem}{thm:p15-sequential-leading-minimax}{proved=annex:p15.sequential.leading.minimax}{Leading cost with finite entropy, and the singular endpoint}
Use \(a=|\log\beta|\), \(b=|\log\alpha|\) in the robust likelihood test.
It satisfies the two robust error constraints.
For \(\alpha,\beta\to0\), with \(t_0=|\log\beta|\),
\(t_1=|\log\alpha|\), if \(D_i<\infty\), then
\[
 \mathcal M_i(\alpha,\beta)\sim \frac{t_i}{D_i},\qquad
 \frac{\sup_{R\in\mathcal H_i}E_RN_{a,b}}
              {\mathcal M_i(\alpha,\beta)}\longrightarrow1.
\]
If \(s_i>0\), then instead
\[
 \mathcal M_i(\alpha,\beta)\longrightarrow1/s_i,\qquad
 \frac{\sup_R E_RN_{a,b}}{\mathcal M_i(\alpha,\beta)}
                 \longrightarrow1.
\]
If \(D_i=\infty\) but \(s_i=0\), then both
\(\mathcal M_i(\alpha,\beta)\) and \(\sup_RE_RN_{a,b}\)
tend to infinity, while their ratios to \(t_i\) tend to zero.
No ratio-optimality assertion between these last two diverging
sublinear quantities is inferred from their zero linear-scale limits.
This distinction does not restrict the full least-favorable ratio
in \cref{thm:p15-robust-endpoint-closure}.
\end{hbtheorem}

\section{Calculated infinite-entropy and singular examples}
\label{sec:p15-endpoint-examples}

\begin{hbtheorem}{thm:p15-infinite-entropy-finite-scores}{proved=annex:p15.infinite.entropy.finite.scores}{Infinite entropy with no infinite observation}
On \(\{0,1,2,\ldots\}\), define
\[
 P(0)=1/2,\quad P(j)=2^{-(j+1)}\ (j\geq1),\qquad
 k=\sum_{j\geq1}2^{-(j+1)}\exp(-2^j),
\]
\[
 Q(0)=1-k,\qquad Q(j)=P(j)\exp(-2^j)\ (j\geq1).
\]
Both laws have full support. Let \(L=\log(2(1-k))\in(0,\log2)\).
Under \(P\), the directional log ratio \(Z=\log(P/Q)\) is
\(-L\) at zero and \(2^j\) at \(j\geq1\). Thus
\(D(P\Vert Q)=\infty\), despite every score being finite.
Write \(m_P(t)=E_P\tau_t\) for the passage mean of these increments.
For \(M=2^m\), integer \(m\geq1\),
\[
 E_P\min(Z,M)=\frac{m+1-L}{2},\qquad
 m_P(2^m)\leq\frac{4\,2^m}{m+1-L}.
\]
The one-sided mean tends to infinity but is sublinear in its boundary.
In the reverse direction,
\[
 D(Q\Vert P)=(1-k)L-\tfrac12\sum_{j\geq1}\exp(-2^j)
                       \in(0,\infty).
\]
With zero neighborhood radii, these are also valid original Huber
endpoint laws. Their two-sided expected stopping time has the
same classification as both boundaries diverge.
\end{hbtheorem}

\begin{hbtheorem}{thm:p15-singular-contamination-stopping}{proved=annex:p15.singular.contamination.stopping}{An exact nonzero-contamination endpoint calculation}
Let \(P_0=\delta_0\), \(P_1=(\delta_0+\delta_1)/2\),
\(\varepsilon_0=0,\varepsilon_1=1/5\).
The least-favorable pair and its likelihood score are
\[
 Q_0=\delta_0,\quad Q_1=(3/5)\delta_0+(2/5)\delta_1,\qquad
 \ell(0)=3/5,\quad\ell(1)=\infty.
\]
Write \(c=\log(5/3)\) and \(r=\lceil a/c\rceil\).
For every finite \(b>0\),
\[
 E_{Q_0}N_{a,b}=r,\qquad
 \sup_{R\in\mathcal H_1}E_RN_{a,b}
 =E_{Q_1}N_{a,b}=\frac{1-(3/5)^r}{2/5}.
\]
The worst null error is zero; the worst alternative error is \((3/5)^r\).
At \(a=3c\), the two expectations are \(3\) and \(49/25\).
Here \(D_0=c<\infty\), \(D_1=\infty\), \(s_1=2/5\);
the alternative expected size tends to \(5/2\), not to zero.
The allowed alternative \(R=P_1\) also has infinite mean directional
score, but its expected stopping time is \(2(1-2^{-r})\), tending to
two. Thus equality of infinite score means does not identify the
score distributions or expected stopping times.

This also shows that errors can tend to zero with the upper boundary
\(b\) fixed. Consequently vanishing errors are not equivalent to
both boundaries diverging in the presence of singular scores.
The general ratio theorem uses the actual error probabilities.
Boundary divergence is an additional sufficient route, and the
infinite-entropy size conclusions use the correct scale.
\end{hbtheorem}

The least-favorable expected-size ratio and the stated leading minimax
results do not assert exact worst expected cost at every pair of finite
boundaries. The latter optimization is a distinct question.
No original contamination class, variation radius convention or
observation support was narrowed to obtain the endpoint closure.
